% !TITLE 忆江南三首
% !AUTHOR 白居易
% !DYNASTY 唐
% !GENRE 词
% !ORDER 50

\begin{poem}
江南好，风景旧曾谙\textsuperscript{1}。
日出江花红胜火，春来江水绿如蓝\textsuperscript{2}。
能不忆江南？

江南忆，最忆是杭州。
山寺月中寻桂子\textsuperscript{3}，郡亭\textsuperscript{4}枕上看潮头\textsuperscript{5}。
何日更重游？

江南忆，其次忆吴宫\textsuperscript{6}。
吴酒一杯春竹叶\textsuperscript{7}，吴娃\textsuperscript{8}双舞醉芙蓉。
早晚\textsuperscript{9}复相逢？
\end{poem}

\begin{annotations}
\notetext{1}{谙：熟悉。旧曾谙，是过去曾经熟悉。一个谙字，写出了白居易对江南的深情——江南的一切他都了如指掌。}
\notetext{2}{蓝：蓝草，一种可以提炼靛青染料的植物。江水绿得像蓝草一样，形容春江之绿浓郁欲滴。}
\notetext{3}{桂子：传说月亮里的桂花树种子会飘落人间。杭州灵隐寺有桂树，白居易曾在月夜去寻访飘落的桂子。}
\notetext{4}{郡亭：杭州官署中的亭子。郡是唐代对州的称呼，白居易曾任杭州刺史，故称郡亭。}
\notetext{5}{潮头：钱塘江的大潮。钱塘江潮以八月最盛，潮头高数丈，是杭州的一大奇观。}
\notetext{6}{吴宫：春秋时吴国的宫殿，此处代指苏州。苏州是吴国故都，白居易曾任苏州刺史。}
\notetext{7}{春竹叶：酒名，因酒色青绿如竹叶而得名。一杯春竹叶，写出了苏州生活的闲适与风雅。}
\notetext{8}{吴娃：吴地的美女。娃是古代对美女的称呼，吴地出美女，故称吴娃。}
\notetext{9}{早晚：什么时候。古汉语中早晚即何时，有期盼之意。早晚复相逢，是明知重逢无望的自我安慰。}
\end{annotations}

\begin{appreciation}
这三首词写于白居易晚年闲居洛阳时。他曾在杭州、苏州做过刺史，江南是他一生中最留恋的地方。谙（ān）是熟悉的意思。

第一首总写江南。江南好，风景旧曾谙——江南好啊，那里的风景我曾经多么熟悉。旧曾谙三个字，有回忆的温柔，也有物是人非的感慨。日出江花红胜火，春来江水绿如蓝——太阳出来，江边的花朵红得像火；春天到来，江水绿得像蓝草。红与绿，火与蓝，色彩浓烈得近乎夸张。但这不是虚构，是记忆的颜色——人回忆美好的事物时，总会不自觉地加深它的饱和度。

第二首专忆杭州。山寺月中寻桂子——杭州灵隐寺有桂树，传说月中桂子会飘落人间。白居易曾在月夜到寺中寻访，明知是传说，却愿意相信。郡亭枕上看潮头——郡亭是官署中的亭子，他躺在亭中枕上看钱塘江潮。一个“枕”字，写出了闲适和惬意。何日更重游——什么时候能再去呢？明知年老体衰，重游无望，还是忍不住问。这种明知故问，是老人特有的温柔。

第三首忆苏州。吴酒一杯春竹叶——春竹叶是酒名，色如竹叶之青。吴娃双舞醉芙蓉——吴地的女子翩翩起舞，像醉了的芙蓉花。早晚复相逢——什么时候能再见到呢？三首词，三个问句，一个比一个轻，一个比一个淡。能不忆江南？是肯定的反问；何日更重游？是渺茫的期盼；早晚复相逢？是近乎自语的呢喃。

白居易写这三首词时，已经六十七岁了。江南在他的记忆里，越来越像一个梦。他反复地问，不是真的期待答案，只是在确认那个梦还在。晚年白居易的诗，往往于平淡中见深致。这三首词没有一字写愁，但字字是愁——对逝去时光的追念，对衰老身体的无奈，对远方故地的眷恋。江南好，只是回不去了。
\end{appreciation}
